
\label{appendix:latent_planning}

In this section, we provide a principled explanation for why latent planning 
is substantially more effective than direct action-space sampling in long-horizon
Vision--Language--Action (VLA) tasks.
Our analysis is not intended as a tight optimality guarantee,
but rather as a geometric and probabilistic characterization of the search difficulty
inherent to high-dimensional trajectory spaces, and how latent planning alleviates it
by reweighting probability mass toward feasible trajectories.

\paragraph{Setup.}
Let the environment be a Markov decision process
$(\mathcal{S}, \mathcal{A}, \mathcal{P}, R, \gamma)$.
A planning horizon-$H$ trajectory lies in the space
\[
\mathcal{X} = \mathcal{S}^H \times \mathcal{A}^H,
\]
whose ambient dimension grows linearly with $H$.
We denote by $\mathcal{M}_{\mathrm{traj}} \subset \mathcal{X}$
the set of trajectories that satisfy physical dynamics, contact constraints,
and semantic task requirements.
Throughout, we treat $\mathcal{M}_{\mathrm{traj}}$ as a thin subset of $\mathcal{X}$
with low intrinsic dimension, consistent with the manifold hypothesis
commonly assumed in robotics and control.

We consider a learned latent trajectory generator
$\mathcal{W}_\theta : \mathbb{R}^{d_z} \rightarrow \mathcal{X}$,
which maps a latent variable $z$ to a full trajectory,
\[
\tau_{t:t+H} = \mathcal{W}_\theta(z), \qquad z \sim f_\theta(s_t),
\]
where $f_\theta(s_t)$ is a state-conditioned latent distribution.
The induced trajectory distribution is the pushforward measure
\[
P_{\mathrm{latent}} = (\mathcal{W}_\theta)_\# f_\theta .
\]

